\documentclass[11pt]{article}
\usepackage[T1]{fontenc}
\usepackage{lmodern}
\usepackage[a4paper,margin=27mm]{geometry}
\usepackage{amsmath,amssymb,amsthm,mathtools}
\usepackage{microtype,enumitem,booktabs}
\usepackage[hidelinks]{hyperref}
\usepackage{fancyhdr}
\pagestyle{fancy}\fancyhf{}
\fancyhead[L]{\small Product-gradient equations}
\fancyhead[R]{\small Research manuscript, 8 September 2026}
\fancyfoot[C]{\thepage}
\setlength{\headheight}{14pt}
\numberwithin{equation}{section}
\newtheorem{theorem}{Theorem}[section]
\newtheorem{lemma}[theorem]{Lemma}
\newtheorem{proposition}[theorem]{Proposition}
\newtheorem{corollary}[theorem]{Corollary}
\theoremstyle{definition}\newtheorem{definition}[theorem]{Definition}
\newtheorem{example}[theorem]{Example}
\theoremstyle{remark}\newtheorem{remark}[theorem]{Remark}
\newcommand{\C}{\mathbb C}
\newcommand{\N}{\mathbb N}
\newcommand{\Z}{\mathbb Z}
\newcommand{\NN}{\mathbb Z_{\geq0}}
\newcommand{\dd}{\,\mathrm d}
\newcommand{\Div}{\operatorname{div}}
\newcommand{\supp}{\operatorname{supp}}
\newcommand{\Mcal}{\mathcal M}
\newcommand{\Ac}{\mathcal A}
\newcommand{\Pc}{\mathcal P}
\DeclareMathOperator{\ord}{ord}
\DeclareMathOperator{\rank}{rank}
\DeclareMathOperator{\diag}{diag}
\title{Automatic finite order and algebraic normal forms\\
for product-gradient equations with\\
polynomial--exponential forcing}
\author{Research manuscript for the \texttt{whzy3185/math} project}
\date{8 September 2026}
\begin{document}
\maketitle
\begin{abstract}
We study entire solutions of products of independent constant-coefficient
first-order directional derivatives with a polynomial--exponential right-hand
side. The problem includes both equations in Remark~3 of F.~L\"u's
arXiv:2605.09585v1. We allow arbitrary positive integral powers of the factors.
For a nonzero polynomial prefactor, no growth assumption is necessary:
every solution has order exactly the degree of the exponent when that exponent
is nonconstant, and every solution is polynomial when it is constant.
The connected components of the mixed-Hessian graph determine blocks with a
common polynomial exponential phase. All solutions are then parametrized by
polynomial factors satisfying explicit twisted closedness identities and are
reconstructed by radial integrals. For fixed input data, existence reduces to
finitely many factor allocations and homogeneous linear systems, followed by
one scalar normalization. We recover the zero-free ridge-function
classification with all constants specified, obtain an explicit bound on the
number of connected solutions modulo additive constants, and prove a sharp
arithmetic existence criterion for monomial prefactors. Analytic proofs, not
computational experiments, establish the classification and nonexistence
statements.
\end{abstract}
\noindent\textbf{Keywords.} Entire solution; nonlinear partial differential equation;
polynomial prefactor; logarithmic derivative; algebraic normal form.\\
\textbf{MSC 2020.} 32A15, 32A22, 35F20.

\section{Introduction and relation to Remark 3}
Let $\Ac=(a_{ij})\in\operatorname{GL}_n(\C)$ and define
\[
 D_i=\sum_{j=1}^n a_{ij}\frac{\partial}{\partial z_j}.
\]
The equations singled out in Remark~3 of L\"u~\cite{Lu2026}, on page~6 of
the uploaded version, are
\begin{equation}\label{eq:remark}
 \prod_{i=1}^n D_i u=e^g,
 \qquad
 \prod_{i=1}^n D_i u=p e^g,
\end{equation}
where $p,g\in\C[z_1,\ldots,z_n]$. The first equation is already accessible
from Theorem~2 of that paper by the stated invertible linear change of
variables. It is not an independent novelty claim here. The second equation
introduces zeros into the derivative factors. Those zeros cannot simply be
ignored: even $u(x,y)=e^{xy}$ has $u_xu_y=xy e^{2xy}$ and is not a function
of one linear form.

We give a necessary-and-sufficient normal form for the polynomial-prefactor
problem and, with essentially the same proof, for the weighted equation
\begin{equation}\label{eq:weighted-original}
 \prod_{i=1}^n (D_i u)^{\mu_i}=p e^g,
 \qquad \mu_i\in\N,\quad \mu_i\geq1.
\end{equation}
The weights are an extension of the source problem; taking all $\mu_i=1$
recovers exactly \eqref{eq:remark}. The normal form contains no unspecified
entire functions: it uses polynomial factorizations, explicit linear
identities, and definite integrals of known entire integrands.

The relevant background includes Chen--Han~\cite{ChenHan2022},
Xu--Liu--Xuan~\cite{XuLiuXuan2025}, and the finite-order result of
L\"u--Ma~\cite{LuMa2026}. L\"u's preprint removes the finite-order
assumption for the zero-free equation. Our proof below does not use its
Theorem~2 as a black box. Instead, it applies the classical
several-variable logarithmic-derivative estimate of Vitter~\cite{Vitter1977};
a recent explicit two-radius version is given by
Han--Liu~\cite{HanLiu2025}.

\paragraph{Scope of the literature comparison.}
A publisher-indexed 2026 paper by Xu--Ding~\cite{XuDing2026} concerns
product-type equations with a generalized exponential term in $\C^3$.
Its indexed description identifies a finite-order main result. Its full
hypotheses and complete classification were not available for comparison
when this manuscript was prepared. Accordingly, we make no claim of
priority over that paper, or of an exhaustive novelty search. The theorem
proofs below are supplied independently. This is a research manuscript,
not a submitted or peer-reviewed article.

\section{The complete normal form}
Set $I=\{1,\ldots,n\}$ and $M=\sum_{i=1}^n\mu_i$. The exact coordinate
change is
\begin{equation}\label{eq:change}
 z=\Ac^T w,\qquad v(w)=u(\Ac^T w),\qquad
 P(w)=p(\Ac^T w),\qquad G(w)=g(\Ac^T w).
\end{equation}
The chain rule gives $v_{w_i}(w)=(D_i u)(\Ac^T w)$. Thus it suffices to solve
\begin{equation}\label{eq:canonical}
 \prod_{i=1}^n v_i^{\mu_i}=P e^G,
 \qquad v_i=\partial_{w_i}v.
\end{equation}
Until Section~\ref{sec:degenerate}, assume $P\not\equiv0$.
For $B\subseteq I$, write $w_B=(w_i)_{i\in B}$ and let $\iota_B w_B$
be the vector obtained by putting zero in coordinates outside $B$.

\begin{definition}\label{def:data}
For a partition $\Pc$ of $I$ into nonempty blocks, put
\begin{equation}\label{eq:phases}
 \begin{gathered}
 G_0=G(0),\qquad \lambda=e^{G_0/M},\qquad
 m_B=\sum_{i\in B}\mu_i,\\
 G_B(w_B)=G(\iota_Bw_B)-G_0,\qquad H_B=G_B/m_B.
 \end{gathered}
\end{equation}
An admissible polynomial datum for $(P,G,\mu)$ consists of this partition
and nonzero polynomials $q_1,\ldots,q_n$ such that
\begin{align}
 G(w)&=G_0+\sum_{B\in\Pc}G_B(w_B),\label{eq:add}\\
 q_i&\in\C[w_B]\quad(i\in B),\qquad
 \prod_{i=1}^n q_i^{\mu_i}=P,\label{eq:factor}\\
 \partial_jq_i+q_i\partial_jH_B
 &=\partial_iq_j+q_j\partial_iH_B
 \quad(i,j\in B).\label{eq:closed}
\end{align}
All these are polynomial identities. In particular $H_B(0)=0$.
\end{definition}

\begin{theorem}[Complete normal form]\label{thm:normal}
Let $P\not\equiv0$ and $G$ be polynomials, and let all $\mu_i$ be positive
integers. The entire solutions of \eqref{eq:canonical} are exactly
\begin{equation}\label{eq:integral}
 v(w)=C+\lambda\sum_{B\in\Pc}\int_0^1
 e^{H_B(tw_B)}\Bigl(\sum_{i\in B}w_iq_i(tw_B)\Bigr)\dd t,
 \qquad C\in\C,
\end{equation}
where $(\Pc,q_1,\ldots,q_n)$ ranges over the admissible data of
Definition~\ref{def:data}. Each integrand is entire in $w$, uniformly on
compact sets for $0\leq t\leq1$, and the resulting gradient is
\begin{equation}\label{eq:gradient}
 v_i=\lambda q_i e^{H_B}\quad(i\in B).
\end{equation}
For every solution, one admissible partition is the partition into
connected components of the graph with edge $ij$ whenever
$v_{ij}\not\equiv0$, $i\ne j$.
\end{theorem}

The integral formula is an explicit quadrature, not an assertion that every
primitive is elementary. Repeated representations are harmless: the
mixed-Hessian components supply a canonical finest additive decomposition,
whereas the parametrization permits redundant partitions.

\begin{theorem}[Automatic growth and the constant-exponent case]\label{thm:growth}
Every entire solution of \eqref{eq:canonical} has finite order, without
any a priori growth restriction. More precisely:
\begin{enumerate}[label=(\roman*),leftmargin=*]
\item If $d=\deg G\geq1$, then $\rho(v)=d$ and
$T(r,v)=O(r^d)$.
\item If $G$ is constant, then $v$ is polynomial and
\begin{equation}\label{eq:degree}
 \deg v\leq 1+\left\lfloor\frac{\deg P}{\min_i\mu_i}\right\rfloor.
\end{equation}
\end{enumerate}
The same assertions hold for $u$ in \eqref{eq:weighted-original}.
In particular, for Remark~3 with a nonzero polynomial $p$, the growth order
is $\deg g$ when $g$ is nonconstant, and $\deg u\leq\deg p+1$ when $g$
is constant.
\end{theorem}

The proofs occupy the next three sections. The key point is that finite
order is derived before polynomial logarithms are asserted.

\section{Analytic and algebraic preliminaries}
Use the standard Nevanlinna characteristic $T(r,f)$ in $\C^n$, normalized
so that for entire $f$, it equals the spherical mean of $\log^+|f|$ up
to a bounded additive term. We use the first fundamental theorem and
\[
 T(r,fg),\ T(r,f+g)\leq T(r,f)+T(r,g)+O(1),\qquad
 T(r,f^a)=aT(r,f)+O(1)\quad(a\in\N).
\]
For a nonzero polynomial $R$, $T(r,R)=O(\log r)$. The standard
logarithmic-derivative estimate is
\begin{equation}\label{eq:LD}
 m\left(r,\frac{\partial_jf}{f}\right)
 =O\bigl(\log(e+T(r,f))+\log r\bigr)
\end{equation}
outside a set of finite linear measure; see
\cite{Vitter1977,HanLiu2025}. The estimate concerns proximity, not the full
characteristic. The distinction is essential when $P$ has zeros.
One may choose a common translated origin at which the finitely many
functions being considered are nonzero. This meets the origin
normalization in the two-radius formulation and does not change polynomial
degrees or growth order. A finite union of exceptional sets still has
finite measure.

\begin{lemma}[Adding the pole term]\label{lem:LD}
If $f$ is nonzero entire and $\Div(f)\leq\Div(R)$ for a nonzero
polynomial $R$, then, whenever $L_j=(\partial_j f)/f$ is not identically
zero,
\[
 T(r,L_j)=O\bigl(\log(e+T(r,f))+\log r\bigr)
\]
outside a set of finite linear measure.
\end{lemma}
\begin{proof}
At a generic smooth point of a zero hypersurface of $f$, logarithmic
differentiation produces at most a simple pole. Hence the pole divisor of
$L_j$ is bounded by the zero divisor of $f$, and therefore by $\Div(R)$.
Consequently $N(r,L_j)=O(\log r)$. Adding this estimate to
\eqref{eq:LD} proves the assertion. Possible codimension-two singular sets
do not contribute divisorial poles.
\end{proof}

\begin{lemma}[Polynomial divisors and exponential factors]\label{lem:factor}
Suppose nonzero entire functions $f_1,\ldots,f_n$ satisfy
$\prod_i f_i^{\mu_i}=P e^G$, with $P\not\equiv0$ polynomial. Then
\[
 f_i=R_i e^{h_i},\qquad R_i\in\C[w]\setminus\{0\},\quad
 h_i\text{ entire},\qquad \prod_i R_i^{\mu_i}=P.
\]
If, in addition, $T(r,f_i)=O(r^d+\log r)$ for an integer $d\geq0$,
then $h_i$ is a polynomial of degree at most $d$.
\end{lemma}
\begin{proof}
Factor $P=c\prod_{\nu=1}^s\pi_\nu^{e_\nu}$ into distinct irreducible
polynomials. The order of $f_i$ on the irreducible hypersurface
$\{\pi_\nu=0\}$ is a nonnegative integer $a_{i\nu}$ and
$\sum_i\mu_i a_{i\nu}=e_\nu$. Here one uses the elementary divisor fact
that the regular locus of an irreducible complex algebraic hypersurface
is connected; vanishing order is constant there. Local holomorphic
factorization, followed by removable extension across its singular locus,
then gives the global divisibility by $\pi_\nu^{a_{i\nu}}$.
Choose nonzero constants $c_i$ with $\prod_i c_i^{\mu_i}=c$, and set
$R_i=c_i\prod_\nu\pi_\nu^{a_{i\nu}}$. The quotient $f_i/R_i$ and its
reciprocal extend holomorphically across the hypersurfaces. They extend
across the remaining codimension-two set by Hartogs' theorem when
$n\geq2$; in dimension one ordinary removable singularities suffice.
Thus $f_i/R_i$ is zero-free entire. Since $\C^n$ is simply connected,
it has an entire logarithm $h_i$.

For the second assertion,
\[
 T(r,e^{h_i})\leq T(r,f_i)+T(r,1/R_i)+O(1)
 =O(r^d+\log r).
\]
The submean/Poisson estimate for the nonnegative subharmonic function
$\log^+|e^{h_i}|$ bounds its maximum on a ball of radius $r$ by a constant
multiple of its spherical mean on the ball of radius $2r$. Therefore
$\max_{\|w\|\leq r}\Re h_i(w)=O(r^d+\log r)$. Borel--Carath\'eodory
on each complex line through the origin bounds $h_i-h_i(0)$ on the
smaller ball by the same quantity, uniformly in the line direction.
Multivariable Cauchy estimates now show that every Taylor coefficient of
total degree greater than $d$ is zero. For $d=0$, the bound $O(\log r)$
forces every positive-degree coefficient to vanish. This proves the claim.
\end{proof}

\begin{lemma}[Two elementary polynomial facts]\label{lem:elementary}
Let $R\ne0$ and $Q$ be polynomials.
\begin{enumerate}[label=(\roman*),leftmargin=*]
\item If $e^Q$ is rational, then $Q$ is constant.
\item If $\partial_jR+R\partial_jQ=0$, then
$\partial_jR=\partial_jQ=0$.
\end{enumerate}
\end{lemma}
\begin{proof}
For (i), restrict to a generic affine complex line on which a purported
nonconstant $Q$ remains nonconstant and a rational representation has
nonzero denominator. A rational function equal to $e^{Q}$ on that line
has no finite poles, hence is polynomial; a zero-free polynomial on
$\C$ is constant. Differentiation then contradicts nonconstant $Q$.
For (ii), if $\partial_jQ\ne0$, the total degree of
$R\partial_jQ$ is at least $\deg R$, whereas that of a nonzero
$\partial_jR$ is at most $\deg R-1$. They cannot cancel. The case
$\partial_jR=0$ yields the same conclusion immediately.
\end{proof}

\begin{lemma}[Radial primitive]\label{lem:radial}
Let $F_i$ be entire functions of $w_B$, $i\in B$, with
$\partial_jF_i=\partial_iF_j$. Then
\[
 V_B(w_B)=\int_0^1\sum_{i\in B}w_iF_i(tw_B)\dd t
\]
is entire and $\partial_j V_B=F_j$.
\end{lemma}
\begin{proof}
Compact-uniform holomorphy justifies differentiating under the integral.
Using closedness, the derivative of the integrand with respect to $w_j$ is
\[
 F_j(tw_B)+t\sum_{i\in B}w_i\partial_jF_i(tw_B)
 =F_j(tw_B)+t\sum_{i\in B}w_i\partial_iF_j(tw_B)
 =\frac{\dd}{\dd t}\bigl[tF_j(tw_B)\bigr].
\]
Integration from $0$ to $1$ gives $F_j(w_B)$. The lower endpoint vanishes.
\end{proof}

\section{The growth bootstrap without a finite-order assumption}
\begin{proposition}\label{prop:bootstrap}
For any entire solution of \eqref{eq:canonical}, set
$d=\deg G$ when $G$ is nonconstant and $d=0$ otherwise. Then
\begin{equation}\label{eq:bootbound}
 T(r,v_i)=O(r^d+\log r)\qquad(1\leq i\leq n).
\end{equation}
\end{proposition}
\begin{proof}
Write $f_i=v_i$. No $f_i$ vanishes identically because $P\not\equiv0$.
Let $\Gamma$ be the graph on $I$ with an edge $ij$ exactly when
$v_{ij}\not\equiv0$. If $B$ is one of its connected components, then
$\partial_k f_i=0$ for $i\in B$ and $k\notin B$. Thus $f_i$ depends
only on $w_B$.

Choose $a\in\C^n$ with $P(a)\ne0$. All $f_i(a)$ are nonzero. Freeze
the coordinates outside $B$ at those of $a$ in \eqref{eq:canonical}.
It follows that
\begin{equation}\label{eq:blockprod}
 F_B(w_B):=\prod_{i\in B}f_i(w_B)^{\mu_i}
 =\frac{P(w_B,a_{I\setminus B})}
 {\prod_{k\notin B}f_k(a)^{\mu_k}}
 e^{G(w_B,a_{I\setminus B})}.
\end{equation}
The polynomial in the numerator is not identically zero, as its value at
$a_B$ is $P(a)\ne0$. Consequently, viewed also as a function on $\C^n$,
\begin{equation}\label{eq:blockgrowth}
 T(r,F_B)=O(r^d+\log r).
\end{equation}
This step does not assume a global factorization of $P$ by blocks.

Put $S(r)=1+\sum_iT(r,f_i)$, adding bounded normalization constants if
necessary so that $S$ is positive and nondecreasing. For an edge $ij$,
set $L_{ij}=(\partial_j f_i)/f_i$. Equality of mixed derivatives yields
\begin{equation}\label{eq:ratio}
 f_i L_{ij}=f_j L_{ji}\not\equiv0,
 \qquad \frac{f_i}{f_j}=\frac{L_{ji}}{L_{ij}}.
\end{equation}
The divisor identity in \eqref{eq:canonical} implies
$\Div(f_i)\leq\Div(P)$. Lemma~\ref{lem:LD} and the first fundamental
theorem therefore give
\begin{equation}\label{eq:edgebound}
 T(r,f_i/f_j)=O(\log(e+S(r))+\log r)
\end{equation}
outside a common exceptional set of finite linear measure. Following any
finite path gives the same bound for $i,j$ in a common component, with a
possibly larger fixed constant.

Fix $k\in B$ and write $m_B=\sum_{i\in B}\mu_i$. The identity
\[
 F_B=f_k^{m_B}\prod_{i\in B}(f_i/f_k)^{\mu_i}
\]
implies
\[
 m_BT(r,f_k)\leq T(r,F_B)+\sum_{i\in B}\mu_iT(r,f_i/f_k)+O(1).
\]
After summing over the finitely many vertices and using
\eqref{eq:blockgrowth}--\eqref{eq:edgebound}, we obtain
\begin{equation}\label{eq:absorb}
 S(r)\leq C(r^d+\log r)+C\log(e+S(r))
\end{equation}
outside that exceptional set. There is no comparison of unrelated
``small functions'' in this argument: every error is controlled by the
single nonnegative quantity $S$.

For $x\geq0$, $C\log(e+x)\leq x/2+C'$ with $C'$ depending only on $C$.
Absorbing this term in \eqref{eq:absorb} proves
$S(r)=O(r^d+\log r)$ off the exceptional set. Its tail has measure less
than $1$ for all sufficiently large lower endpoints, so every sufficiently
large interval $[r,r+1]$ contains a point outside it. Monotonicity of $S$
then yields the same bound for every sufficiently large $r$.
This proves \eqref{eq:bootbound}, including singleton components.
\end{proof}

\section{Proof of the normal form and the exact order}
\begin{proof}[Proof of Theorem~\ref{thm:normal}: necessity]
By Proposition~\ref{prop:bootstrap} and Lemma~\ref{lem:factor},
\begin{equation}\label{eq:polylog}
 f_i=R_i e^{Q_i},\qquad R_i\ne0,\quad Q_i\text{ polynomial}.
\end{equation}
For an edge $ij$ of $\Gamma$, define the nonzero polynomials
\[
 A_{ij}=\partial_jR_i+R_i\partial_jQ_i,\qquad
 A_{ji}=\partial_iR_j+R_j\partial_iQ_j.
\]
Mixed-derivative symmetry gives
\[
 A_{ij}e^{Q_i}=A_{ji}e^{Q_j},\qquad
 e^{Q_i-Q_j}=A_{ji}/A_{ij}.
\]
The exponential is rational, so Lemma~\ref{lem:elementary}(i) shows that $Q_i-Q_j$ is constant. Along paths, all $Q_i$ in a component
$B$ have the same nonconstant part.

For $i\in B$ and $k\notin B$, independence of $f_i$ from $w_k$ gives
$\partial_kR_i+R_i\partial_kQ_i=0$. Part (ii) of that lemma gives
$\partial_kR_i=\partial_kQ_i=0$. Normalize the common phase to vanish
at the origin and absorb the resulting exponential constants into the
polynomial factors. Thus
\begin{equation}\label{eq:normalized-component}
 f_i=r_i(w_B)e^{K_B(w_B)},\qquad K_B(0)=0,\quad r_i\ne0.
\end{equation}
Here $r_i$ denotes a polynomial, not a weight. Substituting into the
product equation shows that
\[
 e^{G-\sum_Bm_BK_B}=\frac{\prod_i r_i^{\mu_i}}{P}
\]
is rational. Its polynomial exponent is constant, and evaluation of that
exponent at $0$ shows that the constant is $G_0$. Hence
\[
 G-G_0=\sum_Bm_BK_B,\qquad
 \prod_i r_i^{\mu_i}=e^{G_0}P.
\]
Restricting to the coordinate subspace of $B$ and using $K_{B'}(0)=0$
for $B'\ne B$ yields $m_BK_B=G_B$. Therefore $K_B=H_B$ as defined in
\eqref{eq:phases}. Set $q_i=\lambda^{-1}r_i$. Then
$\prod_iq_i^{\mu_i}=P$, since $\lambda^M=e^{G_0}$. The remaining
mixed-derivative identities are exactly \eqref{eq:closed}.

We have obtained admissible data and \eqref{eq:gradient}.
Lemma~\ref{lem:radial} supplies the potential in \eqref{eq:integral}.
The difference between it and $v$ has every first derivative zero on
connected $\C^n$ and is consequently an additive constant.
\end{proof}

\begin{proof}[Proof of Theorem~\ref{thm:normal}: sufficiency]
Start with admissible data. For each block put $F_i=q_i e^{H_B}$.
Equation~\eqref{eq:closed} says exactly that
$\partial_jF_i=\partial_iF_j$ within the block. Lemma~\ref{lem:radial}
gives \eqref{eq:gradient}; all cross-block mixed derivatives vanish.
Finally,
\[
 \prod_i v_i^{\mu_i}
 =\lambda^M\Bigl(\prod_iq_i^{\mu_i}\Bigr)
     e^{\sum_Bm_BH_B}
 =e^{G_0}P e^{G-G_0}=P e^G.
\]
The integral is globally entire, so no period or branch condition remains.
\end{proof}

\begin{proof}[Proof of Theorem~\ref{thm:growth}]
The proof of Lemma~\ref{lem:factor} and Proposition~\ref{prop:bootstrap}
give
\[
 \log^+\max_{\|w\|\leq r}|f_i(w)|=O(r^d+\log r).
\]
The radial identity
\[
 v(w)-v(0)=\int_0^1\sum_iw_if_i(tw)\dd t
\]
therefore gives $T(r,v)=O(r^d)$ when $d\geq1$.
For the reverse inequality in the order, the highest homogeneous part
$G_d$ is nonzero and
\[
 T(r,e^G)=r^d\int_{\|\zeta\|=1}(\Re G_d(\zeta))^+\dd\sigma(\zeta)
          +O(r^{d-1}+1).
\]
The integral is positive: phase rotation of a nonzero value of $G_d$
produces positive real part on a nonempty open portion of the sphere.
The first fundamental theorem thus gives
$T(r,Pe^G)=\Theta(r^d)$, since multiplication by $P$ changes this
characteristic by at most $O(\log r)$. Cauchy estimates followed by the
submean estimate give
$T(r,f_i)\leq C_nT(4r,v)+O(\log r)$. Applying the product inequality to
\eqref{eq:canonical} yields $\rho(v)\geq d$. Hence $\rho(v)=d$.

When $G$ is constant, \eqref{eq:bootbound} and
Lemma~\ref{lem:factor} show that all $Q_i$ are constant. Thus every
$f_i$ is polynomial and radial integration makes $v$ polynomial.
If $D=\deg v$, at least one derivative, say $f_j$, has degree $D-1$;
none of the other derivatives is identically zero. Therefore
\[
 \deg P=\sum_i\mu_i\deg f_i\geq\mu_j(D-1)
 \geq(\min_i\mu_i)(D-1),
\]
which is \eqref{eq:degree}. Invertible linear changes preserve total
polynomial degree and entire-function order, proving the assertions for
$u$. Sharpness in the unweighted case is witnessed by
$v=w_1^{D}/D+\sum_{i=2}^n w_i$, whose derivative product is $w_1^{D-1}$.
\end{proof}

\section{Finite algebraic existence testing and a counting bound}
Factor the specified input over $\C$:
\begin{equation}\label{eq:inputfactor}
 P=c_P\prod_{\nu=1}^s\pi_\nu^{e_\nu},\qquad c_P\ne0.
\end{equation}
For every admissible datum,
\begin{equation}\label{eq:allocation}
 q_i=c_iR_i,\qquad
 R_i=\prod_\nu\pi_\nu^{a_{i\nu}},\qquad
 a_{i\nu}\in\NN,\quad \sum_i\mu_i a_{i\nu}=e_\nu,\qquad
 \prod_i c_i^{\mu_i}=c_P.
\end{equation}
There are only finitely many integer allocations. The constants remain
unknown, but their equations are especially simple.

\begin{theorem}[Finite linear-algebra reduction]\label{thm:algorithm}
For fixed $P,G,\mu$, enumerate all partitions of $I$ and the allocations
\eqref{eq:allocation}. Reject a partition if \eqref{eq:add} fails, and
reject an allocation if $R_i$ depends on a variable outside its block.
For each remaining pair form the coefficient matrix of the polynomial
identities
\begin{equation}\label{eq:linear}
 c_i\bigl(\partial_jR_i+R_i\partial_jH_B\bigr)
 -c_j\bigl(\partial_iR_j+R_j\partial_iH_B\bigr)=0
 \quad(i<j,\ i,j\in B).
\end{equation}
Let $L\subseteq\C^n$ be its nullspace. This pair supplies an entire
solution if and only if no coordinate functional is identically zero on
$L$. If it does, choose $c\in L$ with all coordinates nonzero and multiply
it by a scalar $t$ satisfying
\begin{equation}\label{eq:scale}
 t^M=\frac{c_P}{\prod_i c_i^{\mu_i}}.
\end{equation}
The scaled constants give all solutions through Theorem~\ref{thm:normal},
when all such choices are permitted.
\end{theorem}
\begin{proof}
Necessity of the integer allocation follows from unique factorization;
\eqref{eq:linear} is \eqref{eq:closed} with $q_i=c_iR_i$.
A nonzero-coordinate vector cannot exist if some coordinate functional
vanishes on $L$. Conversely, if none vanishes identically, the intersections
of $L$ with the coordinate hyperplanes are finitely many proper linear
subspaces. Their union cannot cover the vector space $L$ over the infinite
field $\C$. Thus $L$ contains $c$ with all coordinates nonzero. All
closedness equations are homogeneous, so $tc\in L$. Equation~\eqref{eq:scale}
has a nonzero complex solution and enforces \eqref{eq:allocation}.
The normal form proves sufficiency and exhaustiveness.
\end{proof}

This is a finite symbolic decision procedure over an effective exact
coefficient field after factorization over its algebraic closure. It is
not a numerical decision algorithm for arbitrary uncomputable complex
constants. No degree search for an unknown entire function is involved.

\begin{corollary}[Connected solutions are finite in number]\label{cor:count}
Define
\[
 N_\mu(e)=\#\{(a_1,\ldots,a_n)\in\NN^n:\ \sum_i\mu_i a_i=e\}.
\]
The number of solutions of \eqref{eq:canonical} with connected
mixed-Hessian graph, modulo addition of constants, is at most
\begin{equation}\label{eq:count}
 M\prod_{\nu=1}^s N_\mu(e_\nu).
\end{equation}
In the unweighted case this bound is
$n\prod_\nu\binom{e_\nu+n-1}{n-1}$.
If $G$ cannot be additively separated over any nontrivial coordinate
partition, or if one irreducible factor of $P$ depends on every coordinate,
then every solution has connected mixed-Hessian graph, and the same
bound applies to all solutions modulo constants.
\end{corollary}
\begin{proof}
For a connected solution, the single phase is fixed:
$H=(G-G_0)/M$. Fix one allocation and write
$A_{ij}=\partial_jR_i+R_i\partial_jH$. In the presence of a vector
$c$ with all entries nonzero, \eqref{eq:linear} forces $A_{ij}=0$ exactly
when $A_{ji}=0$. Along each nonzero edge it determines a unique ratio
$c_i/c_j$. Connectedness therefore determines all $c_i$ up to one common
nonzero scalar. Its normalization has at most $M$ choices. Summing over
allocations proves the bound.

A disconnected graph would give \eqref{eq:add}, contradicting the first
additional hypothesis. Under the second, the full-support irreducible
factor must divide at least one $q_i$, but that $q_i$ depends only on its
proper block. A factor of a polynomial independent of $w_j$ is itself
independent of $w_j$ (compare degrees in $w_j$); this is a contradiction.
\end{proof}

\section{Explicit resolution of the zero-free equation}
\begin{corollary}[Weighted ridge-block form]\label{cor:ridge}
For $P\equiv1$, all entire solutions of \eqref{eq:canonical} are given by
a partition $\Pc$, constants $b_i\in\C^*$, and one-variable polynomials
$h_B$ satisfying
\begin{equation}\label{eq:ridge-data}
 \prod_i b_i^{\mu_i}=1,\qquad h_B(0)=0,\qquad
 G=G_0+\sum_{B\in\Pc}m_B h_B(\ell_B),\qquad
 \ell_B(w_B)=\sum_{i\in B}b_iw_i.
\end{equation}
The solutions themselves are
\begin{equation}\label{eq:ridge-solution}
 v(w)=C+e^{G_0/M}\sum_{B\in\Pc}
          \int_0^{\ell_B(w_B)}e^{h_B(s)}\dd s.
\end{equation}
Conversely every such choice solves the equation.
\end{corollary}
\begin{proof}
Unique factorization and $\prod q_i^{\mu_i}=1$ force $q_i=b_i$ to be
nonzero constants. Equation~\eqref{eq:closed} becomes
$b_i\partial_jH_B=b_j\partial_iH_B$. Thus every constant vector tangent
to $\ker\ell_B$ annihilates $H_B$. Complete $\ell_B$ to an invertible
linear coordinate system; derivatives in all the other directions vanish.
It follows that $H_B=h_B(\ell_B)$ for a polynomial $h_B$ with $h_B(0)=0$.
Its gradient primitive is exactly the integral in
\eqref{eq:ridge-solution}. Conversely, differentiation gives
$v_i=e^{G_0/M}b_i e^{h_B(\ell_B)}$, and \eqref{eq:ridge-data} gives the
required product.
\end{proof}

\paragraph{Equation (1.8) of Remark 3.}
Set every $\mu_i=1$, so that $M=n$ and $m_B=|B|$. In the preceding
corollary take $G(w)=g(\Ac^Tw)$ and substitute $w=\Ac^{-T}z$ in
\eqref{eq:ridge-solution}. This is the full pullback formula for $u$.
Singleton blocks are the isolated set $J$ in the source paper; larger
connected blocks are its equivalence classes under $\approx$. The factor
$e^{g(0)/n}$ and the product condition $\prod b_i=1$ keep every constant
normalization explicit. If $g$ is constant, all $h_B$ vanish and $u$ is
affine, recovering the source's rigidity result.

\paragraph{Equation (1.9) of Remark 3.}
For $p\not\equiv0$, use Theorem~\ref{thm:normal} with all weights one,
$P=p\circ\Ac^T$, $G=g\circ\Ac^T$, and then $u(z)=v(\Ac^{-T}z)$.
This is a necessary-and-sufficient classification, not just a family of
examples or a finite-order conditional result.

\section{A sharp arithmetic obstruction for monomial prefactors}
The next result gives closed formulas and exact nonexistence conditions
for an infinite class of genuinely coupled inputs.

\begin{theorem}[Monomial rigidity]\label{thm:monomial}
Let $n\geq2$, $s(w)=w_1\cdots w_n$, let $h\in\C[t]$ be nonconstant,
and let $k_1,\ldots,k_n\in\NN$. The equation
\begin{equation}\label{eq:monomial}
 \prod_{i=1}^n v_i^{\mu_i}
 =\Bigl(\prod_{j=1}^n w_j^{k_j}\Bigr)e^{M h(s(w))}
\end{equation}
has an entire solution if and only if there is an integer $b\geq1$ such
that
\begin{equation}\label{eq:arith}
 k_j=Mb-\mu_j\qquad(1\leq j\leq n).
\end{equation}
When this condition holds, all solutions are
\begin{equation}\label{eq:monomialsolutions}
 v(w)=C+\zeta\int_0^{s(w)}t^{b-1}e^{h(t)}\dd t,
 \qquad \zeta^M=1,\quad C\in\C.
\end{equation}
For the original unweighted product and $k_1=\cdots=k_n=k$, existence
is equivalent to $k\equiv n-1\pmod n$ (with $k\geq0$). In particular,
for $n=2$ exactly the odd $k$ are admissible.
\end{theorem}
\begin{proof}
Every nonconstant monomial of $h(s)$ involves every coordinate, so
\eqref{eq:add} admits only the one-block partition. Theorem~\ref{thm:normal}
therefore gives $H=h(s)-h(0)$ and $\lambda=e^{h(0)}$. The only irreducible
factors of $P$ are coordinates, hence
\[
 q_i=c_i w^{a_i},\quad a_i=(a_{i1},\ldots,a_{in})\in\NN^n,
 \quad\sum_i\mu_i a_{ij}=k_j,
 \quad\prod_i c_i^{\mu_i}=1.
\]
Let $W(s)=s h'(s)$, a nonconstant polynomial with $W(0)=0$.
On the algebraic torus $(\C^*)^n$, closedness is
\begin{equation}\label{eq:mono-closed}
 c_i w^{a_i-e_j}(a_{ij}+W(s))
 =c_j w^{a_j-e_i}(a_{ji}+W(s))\qquad(i\ne j).
\end{equation}
Neither parenthesized polynomial vanishes identically.
The quotient of the two monomials is a rational function of $s$.
Keeping $s$ fixed and replacing $w_r,w_t$ by $\tau w_r,\tau^{-1}w_t$
shows that all entries of $a_i+e_i-a_j-e_j$ are equal. Thus this vector
is $d_{ij}(1,\ldots,1)$ for an integer $d_{ij}$. Equation
\eqref{eq:mono-closed} becomes
\[
 \frac{c_i}{c_j}s^{d_{ij}}
 =\frac{a_{ji}+W(s)}{a_{ij}+W(s)}.
\]
As $s$ tends to infinity away from finitely many poles, the right-hand
side tends to $1$. It follows that $d_{ij}=0$ and $c_i=c_j$.
The same identity now implies $a_{ij}=a_{ji}$.

Consequently all vectors $a_i+e_i$ coincide, say with $\beta$.
For $i\ne j$, $a_{ij}=\beta_j$ and $a_{ji}=\beta_i$, so all entries
of $\beta$ equal an integer $b$. Since $a_{ii}=b-1\geq0$, we have
$b\geq1$. Therefore $a_i=b(1,\ldots,1)-e_i$, all $c_i=\zeta$, and
$\zeta^M=1$. The weighted column sums give exactly
$k_j=Mb-\mu_j$. Moreover,
\[
 v_i=\zeta s^{b-1}\frac{\partial s}{\partial w_i}e^{h(s)},
\]
so radial integration, or differentiation of a one-variable primitive,
gives \eqref{eq:monomialsolutions} up to a constant.
Conversely those entire functions have derivative product
$\zeta^M e^{Mh(s)}\prod_jw_j^{Mb-\mu_j}$, proving sufficiency.
\end{proof}

\section{Examples and the two-variable specialization}
\begin{example}[A non-ridge exponential]
Take $n=2$, all weights one, $P=xy$, $G=2xy$. Theorem~\ref{thm:monomial}
gives exactly $v=C\pm e^{xy}$. For $e^{xy}$,
\[
 \det\nabla^2v=-(1+2xy)e^{2xy}\not\equiv0.
\]
A function of one linear form has Hessian rank at most one. Thus the
zero-free ridge conclusion cannot be copied unchanged to $P\not\equiv1$.
The inputs $P=(xy)^2$, $G=2xy$ admit no entire solution, by the same
theorem; this is an analytic nonexistence statement.
\end{example}

\begin{example}[Polynomial amplitude]
For $v=(x+y)e^x$,
\[
 v_x=(x+y+1)e^x,\qquad v_y=e^x,\qquad
 v_xv_y=(x+y+1)e^{2x}.
\]
Here $q_1=x+y+1$, $q_2=1$, $H=x$, and the closedness identity is
$\partial_yq_1=1=q_2\partial_xH$. The prefactor is irreducible and
depends on both coordinates, so all solutions for these inputs are
connected. The alternative factor allocation $q_1=c$, $q_2=c^{-1}(x+y+1)$
fails closedness, while the displayed allocation forces the two constants
to agree. Hence the complete family is $C\pm(x+y)e^x$.
\end{example}

\begin{example}[A genuinely non-elementary quadrature]
Let $F'(t)=e^{t^2}$ and $F(0)=0$. Then $v=F(xy)$ satisfies
$v_xv_y=xy e^{2x^2y^2}$. Theorem~\ref{thm:monomial} gives precisely
$C\pm F(xy)$. There is no one-variable polynomial $R$ with
$F=Re^{t^2}+C$, because this would require $R'+2tR=1$, impossible by
comparison of degrees. The general normal form must therefore allow
quadratures rather than only polynomial multiples of exponentials.
\end{example}

\begin{example}[Different phases on different blocks]
In three variables let
\[
 v=e^{xy}+z^2/2,
 \quad P=xyz,\quad G=2xy.
\]
Use the blocks $\{1,2\}$ and $\{3\}$ with
$(q_1,q_2,q_3)=(y,x,z)$ and $(H_{\{1,2\}},H_{\{3\}})=(xy,0)$.
This illustrates why a single common phase for all coordinates would be
too restrictive. The derivative product is $xyz e^{2xy}$.
\end{example}

\begin{example}[Unequal powers]
Take $(\mu_1,\mu_2)=(1,2)$, $M=3$, $h(t)=t$.
For $P=x^2y$ condition \eqref{eq:arith} holds with $b=1$, and the
complete family is $v=C+\zeta e^{xy}$, $\zeta^3=1$.
For $P=xy$ the two requirements in \eqref{eq:arith} are incompatible,
so no entire solution exists.
\end{example}

\paragraph{All two-variable inputs.}
For the unweighted equation $v_xv_y=P e^G$ there are just two partition
types. For $\{\{1\},\{2\}\}$ one must have
\[
 G=G_0+G_1(x)+G_2(y),\quad P=q_1(x)q_2(y),
\]
and
\[
 v=C+e^{G_0/2}\left(\int_0^xq_1(t)e^{G_1(t)}\dd t+
                        \int_0^yq_2(t)e^{G_2(t)}\dd t\right).
\]
For the one-block partition put $H=(G-G_0)/2$. Enumerate $q_1q_2=P$
and impose the single polynomial identity
\[
 \partial_yq_1+q_1H_y=\partial_xq_2+q_2H_x.
\]
The complete solutions are then \eqref{eq:integral}. These two cases
exhaust all entire solutions, including polynomial ones, and overlap
only through redundant representations.

\section{Degenerate inputs and the limits of the hypotheses}\label{sec:degenerate}
\begin{proposition}[The zero prefactor]\label{prop:zero}
If $p\equiv0$, the entire solutions of \eqref{eq:weighted-original} are
exactly those for which at least one $D_i u$ vanishes identically.
Equivalently, after \eqref{eq:change}, $v$ is independent of at least one
coordinate. The other $n-1$ variables may enter through an arbitrary
entire function.
\end{proposition}
\begin{proof}
The ring of entire functions on connected $\C^n$ is an integral domain.
A finite product of positive powers is identically zero only if one of
its factors is identically zero. In canonical coordinates $v_i=0$
is equivalent to independence of $w_i$. The converse is immediate.
\end{proof}

For $n\geq2$, Proposition~\ref{prop:zero} allows infinite-order functions,
so the nonzero-prefactor hypothesis is essential to the growth theorem.
The invertibility of $\Ac$ is also essential: with two identical rows
$(1,0)$ the unweighted equation is $u_x^2=1$, and
$u=x+F(y)$ works for any entire $F$. Finally, polynomiality of $G$ is
not a cosmetic assumption. L\"u's Remark~2 and Example~1 already exhibit
failure of the zero-free ridge classification for a transcendental entire
exponent. No such extension is asserted here.

\section{Verification, provenance, and conclusions}
The accompanying exact symbolic code constructs the matrices in
Theorem~\ref{thm:algorithm}, validates the polynomial identities in a
proposed datum, and checks examples and low-dimensional monomial inputs.
These computations verify finite algebraic instances only. They do not
prove Proposition~\ref{prop:bootstrap}, do not certify the general
normal form, and do not establish novelty. Those roles belong respectively
to the analytic proofs and to a separate literature review. No Lean or
other proof-assistant verification is claimed.

For the two questions actually contained in Remark~3, the outcome is as
follows. Equation (1.8) has the explicit pullback ridge-block form of
Corollary~\ref{cor:ridge}. Equation (1.9), for $p\not\equiv0$, has the
finite polynomial-data classification of Theorem~\ref{thm:normal} and
Theorem~\ref{thm:algorithm}, with automatic finite order and a sharp
constant-exponent degree bound. The case $p\equiv0$ is exactly
Proposition~\ref{prop:zero}. Weighted products, the connected-solution
bound, and the monomial arithmetic criterion are additional extensions.
A full-text comparison with \cite{XuDing2026} and an independent expert
proof audit are required before any submission-level novelty or
publication-readiness claim.

\begin{thebibliography}{9}
\bibitem{ChenHan2022} W.~Chen and Q.~Han,
On entire solutions to eikonal-type equations,
\emph{J. Math. Anal. Appl.} \textbf{506} (2022), no.~1, 124704.
\url{https://doi.org/10.1016/j.jmaa.2020.124704}.
\bibitem{HanLiu2025} Q.~Han and J.~Liu,
A short proof of the lemma of the logarithmic derivative in several complex variables,
\emph{Complex Anal. Oper. Theory} \textbf{19} (2025), Article~92.
\url{https://doi.org/10.1007/s11785-025-01701-x}.
\bibitem{Lu2026} F.~L\"u,
On entire solutions for a class of product-type nonlinear PDEs in $\C^n$,
arXiv:2605.09585v1, 10 May 2026.
\url{https://arxiv.org/abs/2605.09585}.
\bibitem{LuMa2026} F.~L\"u and Z.~Ma,
Entire solutions of product type nonlinear partial differential equations in $\C^n$,
\emph{Glasg. Math. J.} \textbf{68} (2026), no.~2, 257--262.
Published online 12 August 2025.
\url{https://doi.org/10.1017/S0017089525100657}.
\bibitem{Vitter1977} A.~Vitter,
The lemma of the logarithmic derivative in several complex variables,
\emph{Duke Math. J.} \textbf{44} (1977), no.~1, 89--104.
\url{https://doi.org/10.1215/S0012-7094-77-04404-0}.
\bibitem{XuDing2026} H.~Y.~Xu and X.~Ding,
Description of entire solutions of the product type nonlinear PDEs with a
 generalized exponential term in $\C^3$,
\emph{J. Math. Anal. Appl.} \textbf{561} (2026), no.~2, 130680.
\url{https://doi.org/10.1016/j.jmaa.2026.130680}.
Publisher metadata and indexed description checked; full-text comparison pending.
\bibitem{XuLiuXuan2025} H.~Y.~Xu, K.~Liu, and Z.~X.~Xuan,
Results on solutions of several product type nonlinear partial differential
equations in $\C^3$,
\emph{J. Math. Anal. Appl.} \textbf{543} (2025), 128885.
\url{https://doi.org/10.1016/j.jmaa.2024.128885}.
\end{thebibliography}

\end{document}
